\scanpage{55}
\mspage{51}
% archive sheet number: 54
\uuline{\deleted{\unsure{Oblg|Obl|Fig}}At 1}\quad Clair

\uline{\deleted{\unsure{Ob|O}}At 2}\quad Clair \deleted{(c'est l'axiome ensembliste)}

\uline{At 3}\quad Clair

\uline{At 4}\quad Clair.

\deleted{\unsure{Cons|On}} Voyons quel est \unsure{le|l'} \deleted{\unsure{ens des}} \deleted{\unsure{supp}} spatios associé au précédent \unsure{Atelier|atelier}. \deleted{\unsure{Comme}} Soit $\unsure{\mathcal A|\mathscr A}=\mathrm{At}\,\mathrm{spat}(\Sigma)$. \unsure{Considérons|Consid.} \deleted{$(\unsure{\Phi|\varPhi},\rho)$} l'application
\[
\begin{aligned}
\mathcal A&\unsure{\text{\textemdash}|\longrightarrow}\Sigma\\
X=(\Phi,\rho)&\longmapsto\text{plus grand élément de }\Phi,\\
&\hspace{4em}\text{(noté }X^{\unsure{\circ|o}}\text{, cf justification \inserted{plus bas})}
\end{aligned}
\]
Ceci dit, la relation $\mathrel{|\unsure{o|\circ}|}$ dans $\mathcal A$ \unsure{est} induite, via cette application, par la relation de \unsure{$=$|même} nom dans $\Sigma$. On trouve donc une application injective
\[
\Sigma_{\mathcal A}\longrightarrow\Sigma_{\Sigma}
\underset{\inserted{\substack{\text{on vient de}\\\text{le voir}}}}{=}\Sigma.
\]
Mais on a aussi
\[
\begin{aligned}
\Sigma&\longrightarrow\mathcal A\\
\unsure{x|X}&\mathrel{\unsure{\vdash|\mapsto}}\deleted{\unsure{\{}}\bigl(\{x\}\unsure{,}\rho_x\deleted{\unsure{\}}}\bigr)
\qquad \rho_x\ \text{\unsure{la|l'} \inserted{\unsure{Unique|unique}} relation d'ordre sur }\{x\}
\end{aligned}
\]
telle que la relation $\mathrel{|o|}$ sur $\Sigma$ soit induite par celle sur $\mathcal A$, et le composé
\diagram{\begin{tikzcd}[ampersand replacement=\&,column sep=.32em]
\Sigma \& \unsure{\text{\textemdash}|\longrightarrow} \& \mathcal A \& \unsure{\text{\textemdash}|\longrightarrow} \& \Sigma
\end{tikzcd}}
est l'identité.
